\documentclass[11pt,oneside]{article}
\usepackage[utf8]{inputenc}
\usepackage{amsmath, amssymb, amsthm, amsfonts, cite}
\usepackage{geometry}
\geometry{a4paper, margin=1in}

\title{Formal Foundations: Kernel Transitions in Hyperk\"ahler Geometries}
\author{Technical Documentation}
\date{April 2026}

\begin{document}

\maketitle

\section{Algebraic Foundations: The K\"ahler Module}
The derivation module $\Omega_{A/R}$ is defined as the universal object for derivations from the $R$-algebra $A$. Following \cite{grothendieck1967}, we define the module of differentials as the quotient:
\begin{equation}
    \Omega_{A/R} = I / I^2
\end{equation}
where $I$ is the kernel of the multiplication map $\mu: A \otimes_R A \to A$. This ensures that the local "kernel state" is fully captured by the first-order algebraic derivations.

\section{The Geometric Transition}
To evolve the system from algebraic state (the kernel) to geometric structure (the output manifold $M$), we invoke the Calabi-Yau condition. As established by Yau \cite{yau1978}, a compact K\"ahler manifold with $c_1(M) = 0$ admits a unique Ricci-flat metric. 



\section{Hyperk\"ahler Symmetry}
The hyperk\"ahler condition represents the extreme case of the Calabi-Yau requirement, where the holonomy group $\text{Hol}(g)$ is strictly contained in $Sp(n)$. Following the construction by Hitchin et al. \cite{hitchin1987}, the manifold is equipped with a triplet of complex structures $(I, J, K)$ satisfying the quaternion algebra:
\begin{equation}
    I^2 = J^2 = K^2 = IJK = -1
\end{equation}
This triple-symmetry ensures that the manifold remains Ricci-flat ($R_{\mu\nu} = 0$) while maintaining invariant information density across the three fundamental symplectic forms $\omega_I, \omega_J, \omega_K$.



\section{References}
\begin{thebibliography}{9}
\bibitem{grothendieck1967}
Grothendieck, A. (1967). \textit{Éléments de géométrie algébrique: IV. Étude locale des schémas et des morphismes de schémas}. Publications Mathématiques de l'IHÉS.

\bibitem{yau1978}
Yau, S. T. (1978). \textit{On the Ricci curvature of a compact K\"ahler manifold and the complex Monge-Amp\`ere equation, I}. Communications on Pure and Applied Mathematics.

\bibitem{hitchin1987}
Hitchin, N. J., Karlhede, A., Lindstr\"om, U., \& Ro\v{c}ek, M. (1987). \textit{Hyperk\"ahler metrics and supersymmetry}. Communications in Mathematical Physics.

\bibitem{beauville1983}
Beauville, A. (1983). \textit{Variétés Kählériennes dont la première classe de Chern est nulle}. Journal of Differential Geometry.
\end{thebibliography}

\end{document}
